\documentclass[12pt]{book}
\usepackage{precalc1}
\setcounter{chapter}{6}
\begin{document}

\chapter{Logarithms}

A logarithm answers the question the exponential leaves open: given a base $b$ and a
result $y$, what exponent produces it? That exponent is $\log_b y$. Since the
exponential $b^{x}$ takes an exponent to a result and the logarithm takes a result
back to its exponent, the two are inverse functions, and everything about the
logarithm follows from reading the exponential backward.

\section{The Logarithm as an Inverse}

The exponential $f(x)=b^{x}$ with $b>0$, $b\neq1$ is one-to-one: distinct exponents
give distinct results. By the inverse machinery, it has an inverse, and that inverse
is the logarithm base $b$.

\begin{definition}{Logarithm}
For base $b>0$, $b\neq1$, the \emph{logarithm} $\log_b y$ is the exponent to which $b$
must be raised to produce $y$:
\[
\log_b y = x \quad\text{exactly when}\quad b^{x}=y.
\]
\end{definition}

The two statements $b^{x}=y$ and $\log_b y=x$ say the same thing in opposite
directions. Reading one into the other is the single most useful move in the subject.

\begin{keyidea}
$\log_b y$ is an exponent: the exponent that turns $b$ into $y$. The exponential and
logarithm undo each other,
\[
b^{\log_b y}=y \qquad\text{and}\qquad \log_b\!\left(b^{x}\right)=x,
\]
because each reverses what the other did.
\end{keyidea}

\begin{example}[evaluating by asking for the exponent]
$\log_2 8 = 3$, because $2^{3}=8$. $\log_{10}1000=3$, because $10^{3}=1000$.
$\log_5 1=0$, because $5^{0}=1$. $\log_3\tfrac19=-2$, because $3^{-2}=\tfrac19$. Each
value is found by asking what exponent on the base yields the argument.
\end{example}

Because the logarithm is the inverse of the exponential, its graph is the reflection
of the exponential graph across the line $y=x$, with domain and range swapped. The
exponential has domain all reals and range $y>0$; the logarithm therefore has domain
$x>0$ and range all reals. Logarithms of zero and of negatives do not exist, since no
exponent on a positive base produces them.

\begin{center}
\begin{tikzpicture}[scale=0.8]
  \draw[->] (-3,0)--(4,0) node[right]{\scriptsize $x$};
  \draw[->] (0,-3)--(0,4) node[above]{\scriptsize $y$};
  \draw[gray,dashed] (-2.8,-2.8)--(3.8,3.8) node[right]{\scriptsize $y=x$};
  \draw[pcblue,thick,domain=-2.8:1.35,smooth,samples=80] plot (\x,{2^\x});
  \node[pcblue] at (1.55,3.6) {\scriptsize $b^{x}$};
  \draw[pcred,thick,domain=0.13:3.8,smooth,samples=80] plot (\x,{ln(\x)/ln(2)});
  \node[pcred] at (3.4,1.4) {\scriptsize $\log_b x$};
  \fill (0,1) circle (1.4pt) node[above left]{\scriptsize $(0,1)$};
  \fill (1,0) circle (1.4pt) node[below right]{\scriptsize $(1,0)$};
\end{tikzpicture}
\end{center}

The exponential passes through $(0,1)$; reflected, the logarithm passes through
$(1,0)$, recording that $\log_b 1=0$ for every base. The $y$-axis is a vertical
asymptote of the logarithm, the reflection of the exponential's horizontal asymptote.

\begin{example}[two named bases]
Base $10$ gives the \emph{common logarithm}, written $\log x$. Base $e$ gives the
\emph{natural logarithm}, written $\ln x$, the inverse of $e^{x}$ promised in the
previous chapter: $\ln b$ is the exponent that turns $e$ into $b$. So $\ln e=1$ and
$\ln 1=0$.
\end{example}

\demo{In the Function Fit Lab, set the parent function to a logarithm and tune
$A\,f(b(x-h))+k$ to match a noisy point cloud. The logarithmic shape rises quickly
near its vertical asymptote and then flattens, the mirror of the exponential's slow
start and rapid climb; fitting one is reading the inverse relationship in
data.}{https://bobovski66.github.io/Precalculus/demos/function\_fit.html}

\section{The Laws of Logarithms}

Each exponent rule, read backward through the inverse, becomes a logarithm rule.
Multiplying powers adds exponents, so a logarithm turns a product into a sum;
dividing powers subtracts exponents, so a logarithm turns a quotient into a
difference; raising a power to a power multiplies exponents, so a logarithm turns an
exponent into a coefficient.

\begin{keyidea}
\textbf{Laws of logarithms.} For $b>0$, $b\neq1$, and positive $M,N$,
\[
\log_b(MN)=\log_b M+\log_b N,\quad
\log_b\!\frac{M}{N}=\log_b M-\log_b N,\quad
\log_b\!\left(M^{p}\right)=p\,\log_b M.
\]
\end{keyidea}

\begin{visual}
A logarithm converts the scale of multiplication into the scale of addition. Products
become sums, quotients become differences, and powers become products. This is why
logarithms turn the wide range of multiplicative quantities, populations, intensities,
probabilities, into manageable additive scales.
\end{visual}

\begin{example}[expanding and combining]
Expand $\log_b\dfrac{x^{3}y}{z}$:
\[
\log_b\frac{x^{3}y}{z}=\log_b x^{3}+\log_b y-\log_b z=3\log_b x+\log_b y-\log_b z.
\]
Run backward, $2\log x+\log y$ combines to $\log(x^{2}y)$, collapsing a sum of
logarithms into the logarithm of a product.
\end{example}

\begin{pitfall}
$\log_b(M+N)$ has no expansion; the laws apply to products, quotients, and powers,
never to sums. In particular $\log_b(M+N)$ is not $\log_b M+\log_b N$. The right-hand
side equals $\log_b(MN)$, an entirely different quantity.
\end{pitfall}

\section{Change of Base}

A calculator computes logarithms in base $10$ and base $e$, yet any base is needed in
practice. The change-of-base formula expresses a logarithm in one base through
logarithms in another. It follows from the inverse relationship: writing
$x=\log_b N$ means $b^{x}=N$; taking $\log_c$ of both sides and using the power law
gives $x\log_c b=\log_c N$.

\begin{keyidea}
\textbf{Change of base.} For positive bases $b,c\neq1$ and positive $N$,
\[
\log_b N=\frac{\log_c N}{\log_c b}.
\]
Taking $c=10$ or $c=e$ lets any logarithm be computed from common or natural
logarithms.
\end{keyidea}

\begin{example}[computing a base-2 logarithm]
Find $\log_2 10$. By change of base,
\[
\log_2 10=\frac{\ln 10}{\ln 2}=\frac{2.302585}{0.693147}\approx 3.3219.
\]
So $2^{3.3219}\approx 10$, which checks: $10$ sits between $2^{3}=8$ and $2^{4}=16$,
closer to $8$.
\end{example}

\begin{example}[solving an exponential equation]
Solve $3^{x}=20$. Take the logarithm of both sides and apply the power law:
\[
x\log 3=\log 20,\qquad x=\frac{\log 20}{\log 3}=\frac{1.30103}{0.47712}\approx 2.727.
\]
The logarithm extracts the exponent, which is exactly what it was built to do.
\end{example}

\section{Logarithms and Information}

The logarithm base $2$ measures information. Consider a quantity that can take one of
$N$ equally likely values, and ask how many yes/no questions are needed to identify
which value occurred. Each question, answered optimally, halves the number of
remaining possibilities. Starting from $N$ and halving down to $1$ takes $\log_2 N$
questions, and that count is the number of \emph{bits} of information the answer
carries.

\begin{keyidea}
The information content of selecting one outcome from $N$ equally likely outcomes is
$\log_2 N$ bits, the number of yes/no questions needed to pin it down. Equivalently,
an event of probability $p$ carries $\log_2\!\dfrac1p=-\log_2 p$ bits: the rarer the
event, the more information its occurrence conveys.
\end{keyidea}

\begin{example}[bits to identify an outcome]
A fair coin has $N=2$ outcomes, so a flip carries $\log_2 2=1$ bit, a single yes/no
answer. A byte chooses one of $N=256$ states, carrying $\log_2 256=8$ bits, since
$2^{8}=256$. Identifying one card from a deck of $52$ needs
$\log_2 52\approx 5.70$ bits: five questions are not quite enough, six more than
suffice.
\end{example}

\begin{example}[information from probability]
An event of probability $p=\tfrac1{32}$ carries $\log_2 32=5$ bits when it occurs,
because $-\log_2\tfrac1{32}=\log_2 32=5$. A near-certain event of probability
$p=0.99$ carries $-\log_2 0.99\approx0.0145$ bits: its occurrence is almost no
surprise, so it conveys almost no information. Rare events are informative; common
events are not.
\end{example}

\begin{visual}
Information is measured on a logarithmic scale because possibilities multiply while
information adds. Two independent yes/no choices give $2\times2=4$ outcomes but
$1+1=2$ bits; ten such choices give $2^{10}=1024$ outcomes and $10$ bits. The
logarithm is the bridge that turns the multiplying count of possibilities into the
adding count of bits.
\end{visual}

\begin{example}[combining independent sources]
A system has two independent components, one with $8$ states and one with $4$. The
combined system has $8\times4=32$ states, carrying $\log_2 32=5$ bits. The same total
comes from adding the parts: $\log_2 8+\log_2 4=3+2=5$. The product law of logarithms
is precisely the additivity of information across independent sources.
\end{example}

\section{Exercises}

\begin{exercises}
\item Evaluate by asking for the exponent: (a) $\log_3 81$; (b) $\log_{10}0.01$;
(c) $\log_2\tfrac1{16}$; (d) $\log_7 1$; (e) $\ln e^{4}$.

\item Rewrite each exponential statement as a logarithmic one and conversely:
(a) $5^{3}=125$; (b) $\log_2 32=5$; (c) $e^{0}=1$.

\item Sketch $f(x)=2^{x}$ and $f^{-1}(x)=\log_2 x$ on one set of axes, mark
$(0,1)$ and $(1,0)$, and identify the asymptote of each.

\item Expand fully: (a) $\log_b(x^{2}yz^{3})$; (b) $\log\dfrac{\sqrt{x}}{y^{4}}$;
(c) $\ln\dfrac{e^{2}x}{y}$.

\item Combine into a single logarithm: (a) $\log x+3\log y$;
(b) $2\ln x-\tfrac12\ln y$; (c) $\log_b 6-\log_b 2$.

\item Explain why $\log_b(M+N)$ cannot be expanded by the laws, and state what
$\log_b M+\log_b N$ does equal.

\item Use change of base to compute, to four decimals: (a) $\log_2 50$;
(b) $\log_5 100$; (c) $\log_3 7$.

\item Solve for $x$, leaving exact and decimal forms: (a) $2^{x}=15$;
(b) $5^{x}=2$; (c) $e^{x}=10$.

\item How many bits of information are carried by identifying one outcome from
(a) $16$ equally likely outcomes; (b) $1000$ equally likely outcomes;
(c) a single roll of a fair six-sided die? Use change of base where needed.

\item An event has probability $p=\tfrac1{64}$. Find the information in bits its
occurrence conveys, and explain why a rarer event carries more information than a
common one.

\item A password is one of $2^{40}$ equally likely strings. State its information
content in bits, and explain how the product law makes the bits of independent
characters add.

\item Two independent dice, one fair six-sided and one fair four-sided, are rolled.
Find the total number of outcomes, the information in bits, and verify that the bits
of the two dice add to the bits of the combined system.
\end{exercises}

\end{document}
